The exact-real model does not bound the encoding length of a repeatedly
updated rational state. We therefore give a specified rounded algorithm,
prove its guarantees for the actual stored iterates, and account for all
integer lengths. This appendix does not assume a positive activation margin.
Throughout, the constant-output cases have been removed, so
$0<\alpha<1$, $0<r<1/d_v\leq1$, and $U=4r<4$.

\subsection{A dimension-independent perturbation bound}

Fix a valid stage. At the actual stored states $\vxi,\vz\in\gK_r$,
form the ideal $\vy,\vq$ by \cref{eq:det-iteration}. Model a rounded step as
\begin{equation*}
 \begin{aligned}
 \vp_0&=\operatorname{Proj}_{\gK_r}(\vq+\vu),&
 \vp&=\vp_0+\ve_p,&
 \vxi^+&=\chi\vxi+\theta\vp+\ve_x,\\
 |\vu|&\leq\kappa_{\mathrm{raw}}\vomega,&
 -\kappa_{\mathrm{proj}}\vomega&\leq\ve_p\leq\vzero,&
 -\kappa_{\mathrm{pri}}\vomega&\leq\ve_x\leq\vzero.
 \end{aligned}
\end{equation*}
We require $\vp,\vxi^+\geq\vzero$. The downward errors then preserve
box and mass feasibility. The ideal projection $\vp_0$ is a comparison
object; the implementation need not enumerate all of its positive entries.

\begin{lemma}[Two energies under directed rounding]
\label{lem:bounded-two-energies}
The energies $E$ and $B$ in \cref{eq:det-first-energy,eq:det-second-energy}
satisfy, at the actual stored states,
\begin{equation}
 \begin{gathered}
 E^+\leq\chi E+e_{\mathrm{step}},\qquad
 B^+\leq\chi B+e_{\mathrm{step}},\\
 e_{\mathrm{step}}:=2\mu_{\mathrm{c}}\kappa_{\mathrm{raw}}
       +\tfrac52\kappa_{\mathrm{pri}}
       +\tfrac52(\theta+\mu_{\mathrm{c}})\kappa_{\mathrm{proj}}.
 \end{gathered}
 \label{eq:bounded-energy-errors}
\end{equation}
For uniform error bounds put $\Gamma=e_{\mathrm{step}}/\theta$. Starting
from zero, the rounded trajectory obeys
\begin{equation*}
 \begin{aligned}
 J_r(\vxi_k)-J_r(\vxi^*)&\leq\chi^k+\Gamma,\\
 \norm{\mQ(\vxi_k-\vxi^*)}_2^2&\leq18\alpha^2r+4\Gamma.
 \end{aligned}
\end{equation*}
\end{lemma}
\begin{proof}
Symmetry and $|\mQ|\vomega=\vomega$ imply
\begin{equation}
 \vomega^\trans|\mQ\va|\leq\vomega^\trans|\va|,
 \qquad |\mQ\va|\leq\kappa\vomega
       \quad\hbox{if }|\va|\leq\kappa\vomega.
 \label{eq:bounded-weighted-contractions}
\end{equation}
Every vector in $\gK_r$ has weighted mass at most one. In particular,
\cref{lem:det-correction-domain} gives $\vxi^*,\vt\in\gK_r$, with
$\vomega^\trans\vt=m_h/\alpha=m_r\leq1$.
The two comparison vectors therefore satisfy
\[
 |(\vp_0-\vxi^*)^\trans\vu|\leq2\kappa_{\mathrm{raw}},
 \qquad
 |(\vp_0-\vt)^\trans\mQ\vu|\leq2\kappa_{\mathrm{raw}}.
\]
The normal of the projection of $\vq+\vu$ has the required sector
signs in both metrics. Replacing that normal by $\vq-\vp_0$ changes
the comparison estimate by at most $2\mu_{\mathrm{c}}\kappa_{\mathrm{raw}}$.
This follows directly from the squared-distance identity in the proof of
\cref{lem:det-comparison}; no nonexpansiveness in a second metric is used.

Let $\vxi_0^+=\chi\vxi+\theta\vp_0$ and
$\vd=\vxi^+-\vxi_0^+=\theta\ve_p+\ve_x\leq\vzero$.
Put $\kappa_d=\theta\kappa_{\mathrm{proj}}+\kappa_{\mathrm{pri}}$.
The lost weighted mass is at most one, so
$\norm{\vd}_2^2\leq\kappa_d\vomega^\trans|\vd|\leq\kappa_d$.
Also
$\vomega^\trans|\mQ\vxi_0^+-\vh|\leq1+m_h\leq2$.
Expanding either $J_r$ or $\mathcal A_r$, their positive mass penalties
decrease under $\vd\leq\vzero$; the remaining linear increase is at
most $2\kappa_d$, by \cref{eq:bounded-weighted-contractions}.
The quadratic increase is at most $\kappa_d/2$, using
$\mQ\preceq\mI$ and $\norm{\mQ\vd}_2\leq\norm{\vd}_2$.

Similarly $\norm{\ve_p}_2^2\leq\kappa_{\mathrm{proj}}$.
Expanding the Euclidean or $\mQ$ mirror-distance term costs at most
$5\mu_{\mathrm{c}}\kappa_{\mathrm{proj}}/2$, because the absolute weighted mass
of the difference between two nonnegative vectors of weighted mass at most
one is at most two.
Adding these contributions proves \cref{eq:bounded-energy-errors}.
Summing the geometric recurrences proves the objective estimate.
The proof of \cref{lem:det-response} with the additional energy $\Gamma$
gives a bound $8\lambda_r m_h+2\Gamma$ on
$\norm{\mQ\vxi_k-\vh}_2^2$, and hence the stated response bound.
\end{proof}

Uniform error per coordinate alone would permit dimension-dependent error.
The estimates above instead use the mass of the actual feasible vectors
and their actual downward losses. Unexposed zero coordinates are not
assigned fictitious independent rounding errors.

\begin{lemma}[Local work survives rounding]
\label{lem:bounded-kinetic-volume}
If $\theta\kappa_{\mathrm{raw}}\leq\lambda_r/4$ and
$\Gamma\leq\alpha^2r$, the actual rounded kinetic supports satisfy
\begin{equation}
 \sum_{k<K}\vol(\supp\vp_k)\leq360K/r.
 \label{eq:bounded-kinetic-volume}
\end{equation}
\end{lemma}
\begin{proof}
Use the comparison set $\gC_r$ from \cref{eq:det-comparison-margin}.
On a selected coordinate $p_i>\xi_i^*$, also $p_{0,i}>\xi_i^*$.
Its lower normal vanishes; the upper and mass normals and $\ve_p\leq0$
can only reduce the selected flow. The raw error contributes at most
$\nu d_i$, where $\nu=\theta\kappa_{\mathrm{raw}}$.
The telescope from \cref{eq:det-selected-flow} therefore becomes
\begin{equation*}
 \frac{\lambda_r}2 V_{\mathrm{out}}
 \leq\sqrt{(V_{\mathrm{out}}+2K/r)H_2}
          +\nu(V_{\mathrm{out}}+2K/r),
 \qquad H_2\leq K(18\alpha^2r+4\Gamma).
\end{equation*}
Downward primal error affects the response estimate; it introduces no
additional term into the identity for the raw kinetic vector at the next
actual state. With $Y=V_{\mathrm{out}}+2K/r$ and $\nu\leq\lambda_r/4$,
the displayed inequality gives
$Y\leq4K/r+4\sqrt{YH_2}/\lambda_r
 \leq4K/r+Y/2+8H_2/\lambda_r^2$.
Consequently $Y\leq8K/r+16H_2/\lambda_r^2\leq360K/r$.
\end{proof}

\subsection{A rounded implementation and its sparse reporter}

Retain the stage parameters $\theta$ and $0<\delta\leq\lambda_r/2$ from
\cref{alg:deterministic-rppr}. Choose a dyadic grid width
$h_{\mathrm{q}}=1/H$, where $H$ is a power of two,
with
\begin{equation*}
 h_{\mathrm{q}}\leq
 \min\{1/8,\theta\tau/256,\delta/2,\lambda_r/8\},
 \qquad \tau=\alpha\delta^2/8.
\end{equation*}
Use the coarsest such grid at the first stage and refine it by halving as
needed at later stages. The grids are thus nested. Let
\begin{equation*}
 T=1/\theta,\qquad K=Tq_{\mathrm{blk}},\qquad
 q_{\mathrm{blk}}:=\min\{j\in\sN:2^j\geq2/\tau\}.
\end{equation*}
Here $T$ is an integer, and $(1-\theta)^T<1/2$ implies
$\chi^K\leq\tau/2$. This finite schedule uses no stored growing exact
power. Its iteration count is $\mathcal{O}(\log(1/\tau)/\sqrt\alpha)$.

Store densities
\begin{equation*}
 \xi_i/\omega_i=\sigma X_i,\qquad Z_i=z_i/\omega_i,
 \qquad \sigma,X_i,Z_i\in h_{\mathrm{q}}\sN_0.
\end{equation*}
At iteration boundaries $1/2\leq\sigma\leq1$.
Maintain primal neighbor records $L_i$, initially
$L_i=\sum_{j\sim i}X_j$, together with exact neighbor sums of $Z_j$
and the current baseline densities. For notational convenience set
$c=(1-\alpha)/2$ and $q_0=(1+\alpha)/2$ in this appendix.
The normalized response used by the reporter is
\begin{equation}
 \widetilde{\mathscr R}_i=(q_0-\mu_{\mathrm{c}})X_i-cL_i/d_i.
 \label{eq:bounded-approximate-response}
\end{equation}
The source is always computed exactly from the stored baseline:
\begin{equation}
 \frac{h_i}{\omega_i}
 =\frac{\alpha\one_{\{i=v\}}}{d_i}
    -q_0\frac{\overline x_i}{\omega_i}
    +\frac c{d_i}\sum_{j\sim i}\frac{\overline x_j}{\omega_j}.
 \label{eq:bounded-exact-source}
\end{equation}
The seed contribution is zero off its one known coordinate.

Use \cref{eq:det-raw-key} with
$\mathscr R_i$ replaced by \cref{eq:bounded-approximate-response} and
with exact kinetic and source terms. The finite reporter computes the exact
box/mass multiplier $\gamma$ for this represented raw vector. Enumerate
only the \emph{closed} tail
\begin{equation}
 k_i\geq\frac{\lambda_r/\theta+\gamma+h_{\mathrm{q}}}{\sigma}.
 \label{eq:bounded-closed-tail}
\end{equation}
At each emitted vertex set
$p_i/\omega_i=\lfloor p_{0,i}/\omega_i\rfloor_{h_{\mathrm{q}}}$, where
$\lfloor t\rfloor_{h_{\mathrm{q}}}=h_{\mathrm{q}}\lfloor t/h_{\mathrm{q}}\rfloor$.
Because $h_{\mathrm{q}}\leq U$, \cref{eq:bounded-closed-tail} is exactly the
condition for a positive stored projection. Equality must be included.
Enumerating positive ideal coordinates smaller than one grid unit would
incur work absent from \cref{eq:bounded-kinetic-volume}; the closed-tail
reporter avoids that enumeration.

Remove the old source and kinetic exceptions. Set
\begin{equation*}
 \sigma'=\lfloor\chi\sigma\rfloor_{h_{\mathrm{q}}},\qquad
 X_i\gets\left\lfloor
 X_i+\frac{\theta(p_i/\omega_i)}{\sigma'}
 \right\rfloor_{h_{\mathrm{q}}}\quad\hbox{at emitted vertices only}.
\end{equation*}
A changed $X_j$ scatters its \emph{actual} dyadic increment to every
neighbor record $L_i$ with $i\sim j$. Set $\vz\gets\vp$ and form the
exact kinetic neighbor sums by scanning these same emitted rows.
Since $\chi\sigma\geq1/4$ and $h_{\mathrm{q}}\leq1/8$, the scale stays
positive. At the old boundary $X_i\leq2U$. Scale rounding therefore loses
at most $2Uh_{\mathrm{q}}$ in physical primal density, and rounding the update
loses at most another $h_{\mathrm{q}}$.

If $\sigma'<1/2$, perform a scalar rebase on all retained records:
\begin{equation*}
 X_i\gets\lfloor\sigma'X_i\rfloor_{h_{\mathrm{q}}},\qquad
 L_i\gets\lfloor\sigma'L_i\rfloor_{h_{\mathrm{q}}},\qquad
 \sigma\gets1.
\end{equation*}
Both replacements use their pre-rebase values. Do not scatter these $X_i$
changes: the neighbor records have already been scaled. Do not scale the
kinetic or baseline neighbor sums. Rebuild all keys from the resulting
records. If $\sigma'\geq1/2$, set $\sigma\gets\sigma'$ instead and
refresh only touched ordinary keys and the new exceptions. No projection
query occurs during an incomplete update. The optional rebase adds at most
one grid unit of downward primal error. Thus
\begin{equation}
 \kappa_{\mathrm{proj}}=h_{\mathrm{q}},\qquad
 \kappa_{\mathrm{pri}}\leq(2U+2)h_{\mathrm{q}}\leq10h_{\mathrm{q}}.
 \label{eq:bounded-downward-errors}
\end{equation}
Nonnegativity is enforced by construction.

\begin{lemma}[Controlled neighbor-response error]
\label{lem:bounded-neighbor-error}
The stored neighbor records satisfy
\begin{equation}
 \frac1{d_i}\left|L_i-\sum_{j\sim i}X_j\right|\leq4h_{\mathrm{q}},
 \qquad \kappa_{\mathrm{raw}}\leq2h_{\mathrm{q}}/\theta.
 \label{eq:bounded-neighbor-error}
\end{equation}
\end{lemma}
\begin{proof}
Let $e_i^L=L_i-\sum_{j\sim i}X_j$. Actual scattered increments preserve
this difference between rebases, including when records are first exposed.
At a rebase, denote its scalar floor losses by $\ell_j$ and $\ell_i^L$.
Then
\[
 e_i^{L,+}=\sigma'e_i^L+\sum_{j\sim i}\ell_j-\ell_i^L,
 \qquad 0\leq\ell_j,\ell_i^L<h_{\mathrm{q}}.
\]
Because $d_i\geq1$ and $\sigma'<1/2$, an old bound of $4h_{\mathrm{q}}$
yields a new bound at most
$\sigma'4h_{\mathrm{q}}+2h_{\mathrm{q}}\leq4h_{\mathrm{q}}$.
The initial error is zero. The only approximate response term is this
primal neighbor sum; its effect on the raw density is at most
\[
 \frac{|u_i|}{\omega_i}
 \leq\frac{\sigma c|e_i^L|}{\theta(1+\theta)d_i}
 \leq\frac{2h_{\mathrm{q}}}\theta.
\]
Retaining the stage's exposed records ensures that this induction does not
silently discard a previously nonzero neighbor contribution.
\end{proof}

Substituting \cref{eq:bounded-downward-errors,eq:bounded-neighbor-error}
into \cref{eq:bounded-energy-errors}, and using $\theta\leq1/2$, gives
\begin{equation}
 \Gamma\leq29h_{\mathrm{q}}/\theta<\tau/8,
 \qquad \theta\kappa_{\mathrm{raw}}\leq2h_{\mathrm{q}}\leq\lambda_r/4.
 \label{eq:bounded-certified-budgets}
\end{equation}
Also $\tau\leq\alpha^2r$. Hence
\cref{lem:bounded-two-energies,lem:bounded-kinetic-volume} certify a gap
below $\tau$ for the actual candidate and a cumulative volume at most
$360K/r$. These statements allow approximate primal neighbor records;
they never identify them with an exact matrix product.

\paragraph{Charging rebases.}
A scale interval starts at one. Each step decreases scale by at most
$\theta+h_{\mathrm{q}}$, so crossing below $1/2$ takes at least
$1/[2(\theta+h_{\mathrm{q}})]$ steps. Since $h_{\mathrm{q}}\leq\theta$,
there are at most $4\theta K+1=\mathcal{O}(1+\log(1/\tau))$ rebases.
Each is a scalar and tree pass over at most $\mathcal{O}(1+K/r)$ retained records;
it inspects no graph adjacency list. Thus rebases contribute only an
additional logarithmic factor. Ordinary scatters and the final candidate
scan retain their explicit cumulative-volume charges.

\subsection{Integer encodings and a finite projection search}

It remains to show that the reporter's arithmetic does not accumulate a
product of degree denominators or one new denominator per iteration.
Write
\[
 \alpha=A_0/B_\alpha,\qquad\theta=1/T,\qquad
 h_{\mathrm{q}}=1/H,\qquad\sigma=S/H,
\]
and express dyadic densities and neighbor records as integers over $H$.
Let $M_i=\sum_{j\sim i}Z_j$. Multiplying the raw key by its degree gives
\begin{equation*}
 \begin{aligned}
 W_i:=d_i k_i={}&-
   \frac{(q_0-\mu_{\mathrm{c}})d_iX_i-cL_i}{\theta(1+\theta)}\\
 &+\frac1\sigma\left\{
   \chi d_iZ_i-\frac{(q_0-\mu_{\mathrm{c}})d_iZ_i-cM_i}{1+\theta}
            +\frac{d_i h_i/\omega_i}{\theta}\right\}.
 \end{aligned}
\end{equation*}
The degree denominator cancels from every term. By
\cref{eq:bounded-exact-source}, all $W_i$ have denominators dividing
$2B_\alpha HT(T+1)S$; the physical weighted keys $\sigma W_i$ have
one common denominator dividing $2B_\alpha H^2T(T+1)$.
Include the represented denominator of $r$ for the common shift and the
box height. A comparison of two unweighted keys uses their two individual
degrees, whereas subtree weighted sums use this one common denominator.
No least common multiple of encountered degrees is formed.

For a direct implementation, partition the records into two disjoint
balanced trees. Clearing the fixed stage denominators expresses the physical
pre-shift density as
\begin{equation*}
 \frac{S\mathfrak b_i+H\mathfrak e_i}{G_0d_i},
\end{equation*}
where $\mathfrak b_i,\mathfrak e_i$ are integers, $G_0$ is a fixed
positive integer, and $\mathfrak e_i=0$ outside the current source and
kinetic exceptions. The ordinary tree stores $\mathfrak b_i/d_i$ with
global scale $S$; the exception tree stores
$\mathfrak n_i/d_i=(S\mathfrak b_i+H\mathfrak e_i)/d_i$ with scale one.
A registry places each vertex in exactly one tree and records the old key
needed for deletion. Each subtree stores cardinality, degree sum, and
numerator sum. All updates and weighted sums therefore use integers.

A physical clipped-mass query is four tail queries: two trees, each at
the lower and upper clipping thresholds. Its breakpoints, in units scaled
by $G_0$, are the ordered sequences
\[
 \{S\mathfrak b_i/d_i\},\quad\{\mathfrak n_i/d_i\},
 \quad\{S\mathfrak b_i/d_i-G_0U\},
 \quad\{\mathfrak n_i/d_i-G_0U\}.
\]
Let $N$ be the number of retained records. The argument of
\cref{lem:det-reporter} applies to these four lists.
Successively binary search each list by rank, narrowing the bracket until
it has no breakpoint from that list in its interior. Previously excluded
breakpoints remain excluded. There are $\mathcal{O}(\log(N+2))$ queries per list,
each using $\mathcal{O}(\log(N+2))$ integer tree operations. An affine interval
then gives the root by one rational division. Equality returns immediately;
a zero cap returns the zero projection. The reporter still costs
$\mathcal{O}(\log^2(N+2))$ operations plus its stored positive output, independently
of activation margins.

First test the clipped mass at $\gamma=0$, as in \cref{lem:det-reporter};
if it is at most $m_r$, take $\gamma=0$. For the active-mass case,
write $a_i^{\mathrm{den}}$ for a represented raw density after the common
$\lambda_r/\theta$ shift. Let $\gI$ be the free coordinates in the final
interval and $\gU_0$ the saturated coordinates. If $\gI$ is nonempty, the
last division gives
\begin{equation*}
 \gamma=\frac{\sum_{i\in\gI}d_i a_i^{\mathrm{den}}
               +U\sum_{i\in\gU_0}d_i-m_r}
              {\sum_{i\in\gI}d_i}.
\end{equation*}
An empty free set is handled by the exact breakpoint comparison.
The only new divisor is a degree sum. The ideal projected densities are
immediately rounded to the grid, so that divisor never enters the next
stored vector. This is a finite ordered search, not repeated numerical
bisection toward an unknown sign margin.

Normalized primal densities are at most $2U$ at iteration boundaries and
remain bounded by an absolute constant during a rebase step. Neighbor sums
add at most the encoding length of a degree or an exposed-record count.
The grids are nested across stages, so the old baseline is exact on the
new grid. It follows that every integer temporary has length
\begin{equation}
 \mathcal{O}\!\left(L_{\mathrm{in}}+\log H+\log T+
       \log(N+2)+\log(d_{\max}+2)\right).
 \label{eq:bounded-integer-length}
\end{equation}
Here $L_{\mathrm{in}}$ includes the binary lengths of the rational input
numerators and denominators and the maximum length of an encountered vertex
label; $d_{\max}$ is the maximum encountered degree. Intermediate cross products change only the absolute
constant. In particular, there is no linear-in-$K$ denominator growth.

\subsection{Rounded handoff and the bit-complexity theorem}

At completion of a stage, materialize the actual candidate and perform one
\emph{exact} projected-gradient step as in \cref{eq:det-repair-map}.
This scan must use the actual stored primal densities; the approximate
records $L_i$ are not a valid replacement for its matrix product.
The projected-gradient point $\vp$ obeys
$\norm{\vp-\vx_r^*}_2\leq\delta/2$ and has a valid subgradient of
norm at most $\delta/2$, by \cref{eq:det-pg-certificate}. Return the
rounded downward correction
\begin{equation}
 u_i/\omega_i=
    \bigl\lfloor[p_i/\omega_i-\delta]_+\bigr\rfloor_{h_{\mathrm{q}}},
 \qquad \overline\vx^+=\max\{\overline\vx,\vu\}.
 \label{eq:bounded-terminal-repair}
\end{equation}
The total density error is at most
$\delta/2+\delta+h_{\mathrm{q}}\leq2\delta$, and $\vu\leq\vx_r^*$.
On a retained positive coordinate, uniform clipping helps the subsolution
inequality, while the additional grid loss changes the $\mQ$ product by
at most $h_{\mathrm{q}}\omega_i$. Thus
\[
 (\mQ\vu-\vb)_i
 \leq(-\lambda_r+\delta/2-\alpha\delta+h_{\mathrm{q}})\omega_i\leq0.
\]
At a zero coordinate, $(\mQ\vu-\vb)_i\leq0$ follows directly from
the Stieltjes signs. The coordinatewise
maximum preserves safety, as in \cref{lem:det-repair}. Consequently the
actual dyadic baseline satisfies all of \cref{eq:det-repaired-invariant}:
source diffuseness, the support bound, and gap at most $2\delta^2/r$.
It supplies the invariant for the next stage and the requested final gap.
Newly positive neighbors are scanned only after this repair.

\begin{theorem}[Bounded-arithmetic deterministic local acceleration]
\label{thm:bounded-deterministic-rppr}
For rational inputs $\alpha,\rho,\epsobj>0$ in the main problem with
$0<\rho<1/d_v$, the rounded continuation algorithm described in this
appendix returns
$\vzero\leq\widehat\vx\leq\vx_\rho^*$ with objective gap at most
$\epsobj$. Its support volume is at most $1/\rho$. Its fully charged
local operation count is $\softO(1/(\rho\sqrt\alpha))$.

Let $L_{\mathrm{par}}=\log(2/(\alpha\rho\min\{1,\epsobj\}))$ and set
\begin{equation*}
 B:=1+L_{\mathrm{in}}+L_{\mathrm{par}}
       +\log\!\left(2+\frac{L_{\mathrm{par}}}{\rho\sqrt\alpha}\right)
       +\log(d_{\max}+2).
\end{equation*}
All numerical integers have $\mathcal{O}(B)$ bits. With schoolbook integer arithmetic
and exact division, the total bit cost is
$\softO(B^2/(\rho\sqrt\alpha))$. The output stores rational densities
and original degrees. The theorem permits arbitrarily small positive
activation margins and exact threshold equalities. For $\rho\geq1/d_v$,
one degree query identifies the exact zero solution.
\end{theorem}
\begin{proof}
For $\alpha=1$, the exact one-coordinate solution has rational density
$[1/d_v-\rho]_+$ and satisfies the stated bounds directly. Assume
$0<\alpha<1$ for the continuation argument.
The initial baseline and the stage induction are the same as in
\cref{alg:deterministic-rppr}, with \cref{eq:bounded-terminal-repair} at
each handoff. The block schedule, \cref{eq:bounded-certified-budgets}, and
\cref{lem:bounded-two-energies} give the required candidate gap.
The rounded repair closes the induction and the final choice of $\delta$
proves accuracy and support safety.

The rounded cumulative volume, finite reporter, sparse source refreshes,
and scalar rebases have just been charged. Their sum over the halving
schedule uses \cref{eq:det-geometric-work} and is
$\softO(1/(\rho\sqrt\alpha))$. At any stage the number of retained
records is $N=\mathcal{O}(1+L_{\mathrm{par}}/(\rho\sqrt\alpha))$.
Choosing the coarsest allowed nested grids gives
$\log H=\mathcal{O}(L_{\mathrm{par}})$, and $\log T=\mathcal{O}(\log(1/\alpha))$.
Hence \cref{eq:bounded-integer-length} is $\mathcal{O}(B)$.
All comparisons, products, sums, floors, and divisions can be performed
with $\mathcal{O}(B^2)$ bit operations. Label processing is included in $L_{\mathrm{in}}$.
No operation uses an unrepresented real oracle, so multiplication by this
encoding cost proves the bit bound.
\end{proof}
